% !PRINTED 639
\weberpage{639}
%<<CONT
1. \gesperrt{The only roots of unity contained in $\Omega_m$
are the powers of $r$ taken with positive and negative
sign.}

To this attaches a quite general theorem on the roots of
unity, valid, that is, in every algebraic number field, which
we shall now prove.

The theorem reads:

2. \gesperrt{If $\alpha$ is an integer of an algebraic number field
 of degree $n$, and if all the numbers $\alpha_1, \alpha_2, \dots, \alpha_n$
 conjugate to $\alpha$ have the absolute
value 1, then $\alpha$ is a root of unity}\footnote{\gesperrt{Kronecker}, ``Zwei Sätze über Gleichungen mit ganzzahligen
Co\"efficienten''. Crelle's Journal, Bd. 53 (1857). \gesperrt{Minkowski}, ``Geometrie
der Zahlen'', Art. 43.}.

For the proof it is first to be remarked that the absolute
value of the norm of an integer $\omega$ is equal to the product of the
absolute values of the conjugate numbers $\omega_1, \omega_2, \dots, \omega_n$, and,
being an integral rational number, is in any case greater than or
equal to 1. It is therefore not possible that all these absolute
values are less than 1. But if they are all equal to 1, as for the
number $\alpha$ assumed in our theorem, then the norm must be $= \pm 1$,
and $\alpha$ is a unit. If $\beta$ is a second unit which, together with its
conjugates, has the absolute value 1, then the quotient $\alpha : \beta$ has
the same property. If therefore, among the numbers conjugate to
this quotient, even a single one is real, we must have
\[ \frac{\alpha}{\beta} = \pm 1 \]
or $\alpha = \pm \beta$, and this then holds also when $\alpha, \beta$
are replaced by the corresponding numbers of a conjugate
field.

If then $\alpha$ and $\beta$ are neither equal nor opposite, their
ratio is not real, and there accordingly subsists between the
absolute values the inequality:
\[ |\alpha \pm \beta| < |\alpha| + |\beta|. \]
(Cf. the Introduction to the first Band, S. 19.)

We have therefore
\[ |\alpha \pm \beta| < 2, \]
%CONT>>
